\section{Implementation Environment}
\label{sec:implementation-environment}

The proposed SysSecOps process was implemented as a Python-based prototype operating
across a hybrid cloud environment. The environment comprised a public-cloud domain and
a private infrastructure domain connected through a common orchestration and security
assessment workflow. The public-cloud domain provided programmable infrastructure,
deployment, compliance information, and operational telemetry, whereas the private
domain represented Linux-based on-premises systems managed through infrastructure
automation. This arrangement enabled the same operational process and risk evaluation
mechanism described in Chapter~\ref{Chapter3} to be exercised across heterogeneous
infrastructure targets.

The implementation environment was organised into four functional groups: prototype
execution, infrastructure provisioning, security assessment, and operational
monitoring. Table~\ref{tab:implementation-environment} summarises the technologies used
to realise each group. These technologies are implementation choices for evaluating the
proposed process; the process itself is not restricted to these particular products.

\begin{table}[htb!]
\centering
\small
\setlength{\tabcolsep}{4pt}
\caption{Technologies used in the implementation environment.}
\label{tab:implementation-environment}
\begin{tabular}{|L{3.6cm}|L{4.2cm}|L{6.0cm}|}
\hline
\textbf{Environment group} & \textbf{Technology} & \textbf{Role in the prototype} \\
\hline
Prototype execution & Python, command-line interface, and Streamlit & Implemented the
process logic, pipeline entry points, and browser-based operator console. \\
\hline
AI-assisted generation & Amazon Bedrock and OpenRouter & Provided interchangeable
language-model backends for generating and regenerating infrastructure code. \\
\hline
Public infrastructure & AWS CDK and AWS CloudFormation & Defined, synthesized, compared,
and deployed resources in the public-cloud domain. \\
\hline
Private infrastructure & Ansible and Linux nodes & Represented and configured the
on-premises domain of the hybrid environment. \\
\hline
Security assessment & Checkov, \texttt{cfn-lint}, \texttt{ansible-lint}, and a
secret-detection component & Examined infrastructure artifacts for syntax errors,
misconfigurations, insecure practices, and exposed secrets. \\
\hline
Contextual risk inputs & Infracost and AWS Config & Supplied estimated cost information
and the live compliance state used by the risk-scoring mechanism. \\
\hline
Code-risk analysis & Bandit, Semgrep, and a logistic-regression model & Extracted static
analysis features and estimated the probability that generated Python code contained
detector-visible insecurity. \\
\hline
Monitoring and audit & AWS Systems Manager, Amazon CloudWatch, AWS CloudTrail, Amazon
EventBridge, AWS Lambda, and Amazon SNS & Collected operational state, recorded activity,
distributed events, and supported notification and remediation actions. \\
\hline
\end{tabular}
\end{table}

\subsection{Prototype Execution Environment}
\label{subsec:prototype-execution-environment}

Python was used as the principal implementation language for the prototype. The
pipeline was accessible through command-line entry points and a Streamlit operator
console. Both interfaces invoked the same orchestration logic, ensuring that an
operation performed through the graphical interface followed the same generation,
assessment, approval, and deployment sequence as an operation initiated from the
command line. External programs were invoked through a shared process-execution layer,
which captured their return status and output in a consistent form for subsequent
processing.

AI-assisted infrastructure generation was provided through interchangeable model
providers. Amazon Bedrock represented a managed cloud-based model service, while
OpenRouter provided an alternative hosted inference interface. The provider was selected
through configuration without changing the remaining pipeline. Generated artifacts,
synthesis output, scanner findings, risk components, and gate decisions were retained
for each run, allowing an unsuccessful generation attempt to be examined and used as
feedback for a subsequent attempt.

\subsection{Hybrid Infrastructure Environment}
\label{subsec:hybrid-infrastructure-environment}

The public part of the environment was hosted on Amazon Web Services. Infrastructure was
expressed as an AWS Cloud Development Kit application written in Python and synthesized
into AWS CloudFormation templates. The synthesized template, rather than only the
high-level source code, was used as the principal public-cloud assessment artifact. This
allowed the security gate to inspect the resource definitions and configuration values
that would be submitted for deployment. Representative resources supported by the
prototype included networking, compute, storage, and database services within a virtual
private cloud.

The private part of the environment consisted of Linux nodes representing on-premises
infrastructure. Their desired configuration was expressed through Ansible inventories
and playbooks. Before execution, the playbooks were subjected to syntax and static
analysis checks. Where operational monitoring across both domains was required, the
private nodes were registered as managed instances through AWS Systems Manager hybrid
activation. This provided a common management path for public-cloud instances and
private Linux nodes while preserving their distinct deployment mechanisms.

Credentials and other sensitive configuration values were supplied to the prototype
through the execution environment rather than embedded in generated infrastructure code
or command-line examples. Deployment remained conditional on the security-gate result:
an accepting result allowed the relevant provisioning mechanism to proceed, a review
result required explicit approval, and a rejection prevented deployment.

\subsection{Security Assessment and Risk Environment}
\label{subsec:security-risk-environment}

The security assessment environment combined several complementary sources rather than
depending on a single scanner. Checkov and \texttt{cfn-lint} evaluated synthesized
CloudFormation templates, while \texttt{ansible-lint} evaluated the private-side
playbooks. A built-in pattern-based component checked the submitted artifacts for
embedded secrets. Their outputs were converted into the common finding representation
defined by the proposed methodology before deduplication and scoring.

Two additional services supplied operational context to the assessment. Infracost
estimated the financial effect of the proposed infrastructure, and AWS Config provided
the compliance state of resources already operating in the public-cloud environment.
Consequently, the gate decision reflected not only static misconfiguration findings but
also cost and live compliance information.

The code-risk component was implemented as a logistic-regression classifier. Bandit and
Semgrep were used to extract severity, confidence, and finding-count features from
Python files. The trained classifier converted these features into a probability of
detector-visible insecurity, which was then incorporated into the unified risk score.
The persisted model and feature scaler were reused during evaluation so that inference
followed the same feature transformation as model training.

Each external assessment component was accessed through an adapter. An unavailable
scanner, missing credential, or unsupported input was recorded as a component status
instead of terminating the complete workflow. The resulting report therefore contained
both the gate decision and the execution status of its contributing assessment sources,
making the conditions under which each result was obtained explicit.

\subsection{Monitoring and Audit Environment}
\label{subsec:monitoring-audit-environment}

After an accepted deployment, the environment used AWS Systems Manager to obtain managed
instance information and support operational actions across the hybrid targets. Amazon
CloudWatch collected gate metrics and operational observations, while AWS CloudTrail
recorded public-cloud API activity for audit purposes. Amazon EventBridge distributed
compliance and deployment events to the operational workflow. Event-driven processing
and notification were supported through AWS Lambda and Amazon SNS, respectively.

The monitoring environment consumed the same target identifiers and gate records
produced during assessment and deployment. This maintained traceability between a
proposed infrastructure change, its security decision, the resulting deployment, and
its subsequent operational state. It also provided the implementation basis for
evaluating whether the proposed process could integrate pre-deployment assessment with
post-deployment monitoring across the public and private domains.
